\section{为什么选择 Diplomacy}

如果只想再做一个超人游戏 AI，围棋和扑克已经提供了成熟舞台。Cicero 团队选择 Diplomacy，是因为现成路线无法简单 scale in：它既有组合式战略行动，又允许玩家用自然语言协商；收益既竞争又合作；真正重要的信息常常存在于承诺、信任和他人预期之中。

\subsection{系统团队本身就是技术证据}

Cicero 不是单一论文技巧的产物。团队 slide 同时列出 reinforcement learning、game theory、NLP、工程和产品化背景的研究者。保留这张看似“行政”的图，是因为它解释了后面的模块化架构：当任务横跨策略与语言时，评估、数据、规划、生成和线上部署必须共同闭环。

\lecturefigure{slide-09-fair-diplomacy-team.jpg}{FAIR Diplomacy 团队：跨 game AI、RL 与 NLP 的系统工程}{官方录像恢复幻灯片 V009；Stanford CS25 Lecture 14}

\teachervoice{老师在此暂停并逐项感谢团队，不只是礼貌。他强调 Cicero 的难点不是找到一个“神奇 loss”，而是让多条研究线在真实多人环境中稳定协作。对复杂 agent 项目，这种 integration work 本身就是核心研究。}

\subsection{两条历史线终于相交}

游戏 AI 的历史从 chess、Go、poker 不断攻克更大的决策空间；语言模型的历史则从 sequence modeling 走向能够生成连贯文本。第十张 slide 把多种游戏里程碑放在一起，第十一张 slide 则用 GPT-2 的独角兽故事提醒观众：到 2019 年，语言生成已经能维持局部连贯，但“会续写”并不等于“会围绕共享世界状态谈判”。

\lecturefigure{slide-10-game-ai-history.jpg}{游戏 AI 里程碑：从棋盘搜索到复杂多人环境}{官方录像恢复幻灯片 V010；Stanford CS25 Lecture 14}

读图时要比较每个任务新增的困难，而不是只看年份：chess 和 Go 主要是完全信息规划；poker 增加隐藏信息与欺骗；实时游戏增加长动作空间与部分可观测性；Diplomacy 再增加自由语言和人类社会约定。每一步都让“状态”更难写成一个封闭张量。

\lecturefigure{slide-11-gpt2-unicorn.jpg}{GPT-2 的流畅生成：语言能力进步与 grounded planning 之间仍有缺口}{官方录像恢复幻灯片 V011；Stanford CS25 Lecture 14}

\begin{warningbox}{语言流畅度不是战略能力代理}
GPT-2 示例证明模型能生成风格连贯文本，却没有证明文本与可执行计划一致。Diplomacy 中一句“我会支持你进入某地”必须对应合法行动、当前单位位置、对手可能反应和未来关系；任何一层脱节，语言越自然反而越容易误导人。
\end{warningbox}

\subsection{规则：同时行动与无约束谈判}

上一节已经说明语言模型与游戏 AI 两条历史线为何要相交；现在必须先把环境的状态转移讲清，否则后面的“grounded”只会成为口号。Diplomacy 由七个国家在欧洲地图上争夺 supply centers。每回合先进行私下或公开沟通，再由所有玩家同时提交行动；承诺不具强制力，玩家可以攻击、支援、运输或按兵不动。由于行动同时揭晓，语言并不是赛后解释，而是会改变联合行动分布的前置控制信号。

\lecturefigure{slide-12-what-is-diplomacy.jpg}{What is Diplomacy：从名称进入规则与任务定义}{官方录像恢复幻灯片 V012；Stanford CS25 Lecture 14}

\lecturefigure{slide-13-diplomacy-rules-challenge.jpg}{七人地图、同步行动、自然语言与超大联合动作空间}{官方录像恢复幻灯片 V013；Stanford CS25 Lecture 14}

\begin{knowledgebox}{联合动作空间为何迅速爆炸}
若第 $i$ 位玩家有合法动作集合 $A_i(s)$，同一回合的联合动作数为 $\prod_{i=1}^{7}|A_i(s)|$。这还没有计入对话历史和隐藏意图。系统不可能简单枚举所有消息与行动组合，因此必须学习其他玩家的 policy，并把 search 集中在高概率、高价值的局部区域。
\end{knowledgebox}

第十四张 slide 把地图与真实聊天并排放置。读这张图应先在地图上找单位和争夺区域，再读消息中提出的合作方案，最后问该承诺会怎样改变双方下单。语言只有落回 board state 和 possible orders，才成为 grounded dialogue。

\lecturefigure{slide-14-diplomacy-map-dialogue.jpg}{Diplomacy 的核心输入：同一个地图状态与不断变化的对话历史}{官方录像恢复幻灯片 V014；Stanford CS25 Lecture 14}

\subsection{Support：为什么沟通会直接改变棋盘结果}

Diplomacy 的单位通常以相同基础力量行动。若一个单位支持另一个单位进攻或防守，联合力量才可能突破。slide 15 用三个小图展示失败、单方 support 成功以及对抗 support 后的结果。这里的教学重点是：跨玩家支援不能由单个玩家独立完成，因此谈判不是装饰，而是合法高价值行动的必要条件。

\lecturefigure{slide-15-support-is-key.jpg}{Support is key：协同承诺如何转化为棋盘上的力量优势}{官方录像恢复幻灯片 V015；Stanford CS25 Lecture 14}

\begin{importantbox}{从消息到状态转移的因果链}
消息提出联合意图；接收者判断可信度并改变 policy；双方提交相互兼容的 orders；规则引擎解析 support；地图状态才发生变化。Cicero 必须建模整条链，而不是只优化下一句最像人类的话。
\end{importantbox}

\subsection{“破坏友谊”与专家真正重视的信任}

Support 规则说明合作能立即改变棋盘，但还没有回答玩家为什么愿意合作。本节因此转向重复互动中的信任成本。媒体常把 Diplomacy 描绘成鼓励背叛的游戏；第十六张 slide 用夸张标题呈现这种印象，第十七张则转向经验丰富玩家的说法：长期稳定合作、可信承诺和可预测性往往比频繁撒谎更有价值。背叛存在，但它会消耗未来合作资本，且这种损失可能跨越多个回合。

\lecturefigure{slide-16-diplomacy-headlines.jpg}{大众叙事：Diplomacy 被称为“会破坏友谊的游戏”}{官方录像恢复幻灯片 V016；Stanford CS25 Lecture 14}

\lecturefigure{slide-17-trust-quote.jpg}{专家视角：在允许背叛的环境里，信任反而更稀缺、更重要}{官方录像恢复幻灯片 V017；Stanford CS25 Lecture 14}

\teachervoice{讲者提醒观众不要把“能骗人”误当作高级智能的充分条件。强玩家会考虑一次谎言之后别人是否还愿意合作；在重复互动里，truthfulness 往往是有战略价值的资源，而不是单纯道德姿态。}

\subsection{为什么它同时是 multi-agent 与 NLP 问题}

规则与信任共同说明，单独优化棋力或单独优化语言都不够。第十八张 slide 用两个相交圆概括任务：reinforcement learning / planning 负责在组合动作空间中做决策，natural language 负责影响其他玩家并建立协作。交集不是把两个模型并排放置，而是让语言产生可测量的策略后果，并让策略模型能判断哪些消息值得相信。

\lecturefigure{slide-18-why-diplomacy.jpg}{为什么选择 Diplomacy：策略学习与自然语言的交集}{官方录像恢复幻灯片 V018；Stanford CS25 Lecture 14}

从 multi-agent 视角看，纯竞争算法会忽略共同收益和人类约定；纯合作算法又可能假设目标完全一致。Diplomacy 处在 mixed-motive 区域：今天的盟友可能是明天的对手，玩家必须推断别人相信什么、期待什么，以及哪些均衡是人类社群中可接受的。

\lecturefigure{slide-19-multi-agent-perspective.jpg}{Multi-agent perspective：混合合作、竞争、常识与人类约定}{官方录像恢复幻灯片 V019；Stanford CS25 Lecture 14}

从 NLP 视角看，已有模型擅长模仿文本，却未必把句子 grounding 到计划。Cicero 追求的是 intentional dialogue：每条消息有明确的沟通目标，能与当前 state、planned action 和收件人的可能反应对齐。

\lecturefigure{slide-20-nlp-perspective.jpg}{NLP perspective：从 imitation 走向 grounded、intentional dialogue}{官方录像恢复幻灯片 V020；Stanford CS25 Lecture 14}

\begin{knowledgebox}{两个视角的共同接口：belief over policies}
规划器需要一个关于“别人接下来会怎么做”的概率模型；对话模块则需要知道“我希望别人怎么做”。这两者都可以落到 policy distribution。语言的作用是改变该分布，search 的作用是在改变后的分布下选择高价值行动。
\end{knowledgebox}

\subsection{本章小结}

Diplomacy 的难点不是规则本身，而是联合动作、自由对话、人类约定和长期信任相互耦合。它迫使系统回答一个比“生成下一句”更严格的问题：这句话是否与可执行计划一致，并且是否会让未来状态更好。

